\section{Lecture 6: Complements on Solid Abelian Groups}
\label{lec:6}

Scholze's last lecture before Clausen resumes is a collection of complements to
the light theory of solid abelian groups built in \cref{lec:5}. It has four
movements. First, the whole picture is upgraded to the derived category: one
defines solidity of a complex, gets a fully faithful $D(\Solid)\hookrightarrow
D(\CondAbl)$ with a derived solidification left adjoint, and observes that
\emph{derived} solidification is where genuine homotopy theory enters --- for a
CW complex $X$ the derived solidification of $\mathbb{Z}[\underline{X}]$ is the
singular chain complex. Second, the abelian category $\Solid$ is analyzed: its
finitely presented objects form an abelian category and are classified, the
crux being a set-theoretic ``freeness'' lemma about finitely generated
submodules of $\prod_{\mathbb{N}}\mathbb{Z}$, and $\prod_{\mathbb{N}}\mathbb{Z}$
turns out to be flat --- a fact special to the light setting. Third, a family of
solid tensor computations is proved, showing that $\otimes^{\solid}$ reproduces
the completed tensor products of formal geometry. Fourth, the same machinery is
run on $p$-adic functional analysis (Smith spaces, Banach spaces, Fréchet
spaces), where the solid tensor product recovers the classical completed tensor
products --- with one derived-duality question that, like the phenomenon of
\cref{lec:4}, depends on the model of set theory.

Throughout, the presentation continues to follow Scholze's first course
\cite{clausen2019condensed}, though the light-specific arguments are new.

\subsection{Recollection}

Recall from \cref{lec:5} the object
\begin{equation}
\label{eq:6:P-recall}
P:=\mathbb{Z}[\mathbb{N}\cup\{\infty\}]\big/\mathbb{Z}[\infty]\in\CondAbl,
\end{equation}
the free light condensed abelian group on a null sequence
(\cref{not:5:P}), which is \emph{internally projective}
(\cref{thm:3:three-properties}(3)), together with the endomorphism
$f=\mathrm{id}-\mathrm{shift}\colon P\to P$, $[n]\mapsto[n]-[n+1]$
(\cref{eq:5:shift}). An object $M\in\CondAbl$ is solid
(\cref{def:5:solid}) iff $f^{\ast}\colon\iHom(P,M)\to\iHom(P,M)$ is an
isomorphism, i.e.\ iff every null sequence in $M$ is uniquely summable. The main
theorem of \cref{lec:5} (\cref{thm:5:structure}, \cref{cor:5:solidification})
was:

\begin{theorem}[Recalled from \cref{lec:5}]
\label{thm:6:recall-structure}
$\Solid\subset\CondAbl$ is an abelian subcategory, stable under all limits and
colimits; the inclusion acquires a colimit-preserving left adjoint
$M\mapsto M^{\solid}$ (\emph{solidification}) and a unique symmetric monoidal
$\otimes^{\solid}$ making solidification symmetric monoidal. It is generated by a
single compact projective object $\prod_{\mathbb{N}}\mathbb{Z}$, with
\begin{equation}
\label{eq:6:generator-tensor-recall}
\prod_{\mathbb{N}}\mathbb{Z}\ \otimes^{\solid}\ \prod_{\mathbb{N}}\mathbb{Z}
\ \cong\ \prod_{\mathbb{N}\times\mathbb{N}}\mathbb{Z},
\end{equation}
and for $S=\varprojlim_n S_n\in\ProFin_{\mathbb{N}}(\mathrm{Fin})$,
\begin{equation}
\label{eq:6:ZS-solid-recall}
\mathbb{Z}[S]^{\solid}\ \xrightarrow{\ \sim\ }\ \varprojlim_{n}\mathbb{Z}[S_{n}].
\end{equation}
\end{theorem}

\subsection{The derived enhancement}

The definition of solidity makes sense verbatim in the derived category. Write
$D(\CondAbl)$ for the (unbounded) derived category of light condensed abelian
groups.

\begin{definition}[Solid complexes]
\label{def:6:derived-solid}
An object $A\in D(\CondAbl)$ is \emph{solid} if
\begin{equation}
\label{eq:6:derived-solid}
f^{\ast}\colon R\iHom(P,A)\ \xrightarrow{\ \sim\ }\ R\iHom(P,A)
\end{equation}
is an isomorphism.
\end{definition}

Because $P$ is internally projective, $R\iHom(P,-)$ is exact, so
$H^{i}\bigl(R\iHom(P,A)\bigr)=\iHom\bigl(P,H^{i}(A)\bigr)$ and
\eqref{eq:6:derived-solid} is equivalent to asking that $f^{\ast}$ be an
isomorphism on $\iHom(P,H^{i}(A))$ for every $i$, that is:
\begin{equation}
\label{eq:6:derived-solid-cohomology}
A\text{ solid}\quad\Longleftrightarrow\quad H^{i}(A)\in\Solid\ \text{ for all }i\in\mathbb{Z}.
\end{equation}
So solidity of a complex may be tested on cohomology groups. This is again
immediate from the internal projectivity of $P$; that completeness for an
analytic ring structure can be checked on cohomology is in general a subtle fact.

\begin{corollary}
\label{cor:6:triangulated}
The class of solid objects of $D(\CondAbl)$ is a triangulated subcategory, stable
under all direct sums, all products and all limits and colimits, and stable
under internal $R\iHom$.
\end{corollary}

The proofs are the same as in \cref{lec:5} (via the exactness of
$R\iHom(P,-)$). Passing to derived categories functorially, one gets:

\begin{proposition}[Derived solidification]
\label{prop:6:derived-solidification}
The functor
\begin{equation}
\label{eq:6:D-solid-embed}
D(\Solid)\longrightarrow D(\CondAbl)
\end{equation}
is fully faithful with essential image the solid objects. The inclusion admits a
colimit-preserving left adjoint, the \emph{derived solidification}
\begin{equation}
\label{eq:6:derived-solidification}
(-)^{\mathbb{L}\solid}\colon D(\CondAbl)\longrightarrow D(\Solid),\qquad
A\longmapsto A^{\mathbb{L}\solid},
\end{equation}
and $D(\Solid)$ carries a unique symmetric monoidal
$\otimes^{\mathbb{L}\solid}$ making $A\mapsto A^{\mathbb{L}\solid}$ symmetric
monoidal; it is the left derived functor of $\otimes^{\solid}$ (and the derived
tensor product is likewise obtained by deriving the one already present on
$\CondAbl$).
\end{proposition}

\begin{remark}[On the existence of the adjoint]
\label{rmk:6:adjoint-existence}
Full faithfulness of \eqref{eq:6:D-solid-embed} is not automatic: it uses that
the compact projective generators stay $\Ext$-vanishing after passage to
$\CondAbl$, i.e.\ that derived solidification of a complex whose terms are direct
sums of the generator $\prod_{\mathbb{N}}\mathbb{Z}$ is computed term by term.
Existence of the left adjoint is the adjoint functor theorem for presentable
$\infty$-categories (Lurie \cite{lurie2017higher}); Scholze noted that theorems
of Neeman and others would also suffice, though they often assume genuine
compact generation, which does not quite hold here, whereas Lurie's version is
definitely sufficient. In any case the adjoint can be constructed by hand, since
it is already identified on the generators. The abstract theorem applies in both
variances: a colimit-preserving functor admits a right adjoint, while a functor
preserving all limits and all sufficiently filtered colimits (i.e.\ an accessible
one) admits a left adjoint.
\end{remark}

On the generators, derived solidification agrees with the underived one --- there
are no higher homotopies:
\begin{equation}
\label{eq:6:derived-solid-generators}
\mathbb{Z}[S]^{\mathbb{L}\solid}\xrightarrow{\ \sim\ }\varprojlim_{n}\mathbb{Z}[S_{n}]
=\mathbb{Z}[S]^{\solid},\qquad
P^{\mathbb{L}\solid}\xrightarrow{\ \sim\ }\prod_{\mathbb{N}}\mathbb{Z},
\end{equation}
and correspondingly
$\prod_{\mathbb{N}}\mathbb{Z}\otimes^{\mathbb{L}\solid}\prod_{\mathbb{N}}\mathbb{Z}
\xrightarrow{\sim}\prod_{\mathbb{N}\times\mathbb{N}}\mathbb{Z}$ is the underived
product.

\subsection{Derived solidification and singular homology}

One might think nothing derived ever appears in practice. It does.

\begin{proposition}[Solidification computes singular homology]
\label{prop:6:singular-homology}
Let $X$ be a CW complex. Then there is a quasi-isomorphism
\begin{equation}
\label{eq:6:singular-homology}
C^{\mathrm{sing}}_{\bullet}(X,\mathbb{Z})\ \xrightarrow{\ \sim\ }\ \mathbb{Z}[\underline{X}]^{\mathbb{L}\solid},
\end{equation}
in particular
$H_{i}^{\mathrm{sing}}(X,\mathbb{Z})\cong H_{i}\bigl(\mathbb{Z}[\underline{X}]^{\mathbb{L}\solid}\bigr)$,
and the derived solidification is a complex of \emph{discrete} abelian groups.
\end{proposition}

This is the homological counterpart of the cohomological statement
$\Ext^{i}(\mathbb{Z}[\underline{X}],\mathbb{Z})\cong H^{i}_{\mathrm{sing}}(X,\mathbb{Z})$
of \cref{thm:4:singular-recovery}, and is again a form of the Dold--Thom
philosophy: $\mathbb{Z}[\underline{X}]$ models the homology of $X$.

\begin{example}[Contractible vs.\ non-contractible]
\label{ex:6:solid-of-spaces}
The free light condensed abelian group $\mathbb{Z}[\underline{[0,1]}]$ on the
interval is a large object (its underlying group is finite sums of points), but
solidifying kills the interval:
\begin{equation}
\label{eq:6:interval-collapse}
\mathbb{Z}[\underline{[0,1]}]^{\mathbb{L}\solid}\simeq\mathbb{Z}.
\end{equation}
For the circle,
\begin{equation}
\label{eq:6:circle}
\mathbb{Z}[\underline{S^{1}}]^{\mathbb{L}\solid}\simeq\mathbb{Z}\oplus\mathbb{Z}[1].
\end{equation}
For a space with homology in unboundedly many degrees --- e.g.\ an infinite
wedge of spheres of all dimensions --- the object $\mathbb{Z}[\underline{X}]$ sits
in degree $0$, yet its derived solidification has homology in arbitrarily high
degrees, i.e.\ is genuinely unbounded.
\end{example}

\begin{remark}[Solidification as a partial passage to homotopy types]
\label{rmk:6:homotopy-types}
Thus the passage $\CondAbl\to\Solid$ is, at least on CW complexes, like passing
to homotopy types: an interval is contracted, and homotopy equivalent CW
complexes are identified. On totally disconnected spaces it retains much finer
information. One inverts homotopy equivalences but \emph{not} weak homotopy
equivalences. The same holds, Scholze noted, for locally contractible spaces (in
the ``funny'' sense of \cref{rmk:4:locally-contractible}); and Mohan Ramachandran
pointed out to him that the paracompactness hypotheses in comparing singular and
sheaf cohomology are actually not needed --- there are recent papers to that
effect.\todo{The ``recent papers'' removing the paracompactness hypothesis
(transcript 00:27:25) may not be by Ramachandran himself; reference not yet
located.}
\end{remark}

\begin{proof}[Proof sketch of \cref{prop:6:singular-homology}]
A formal argument reduces to $X$ a finite CW complex (a filtered colimit in $X$
becomes a limit, and one keeps the finite generation of the homology in each
degree). Assume then $X$ finite CW, hence compact. Choose a surjection
$S\twoheadrightarrow X$ from a light profinite set $S\in\ProFin_{\mathbb{N}}(\mathrm{Fin})$
and form the \v{C}ech nerve; as in \cref{lec:4} its fibre products $S\times_X S,\dots$
are again totally disconnected, giving a resolution
\begin{equation}
\label{eq:6:cech-resolution}
\cdots\to\mathbb{Z}[S\times_{X}S]\to\mathbb{Z}[S]\to\mathbb{Z}[\underline{X}]\to0.
\end{equation}
Applying derived solidification term by term, each $\mathbb{Z}[S_i]^{\mathbb{L}\solid}$
is concentrated in degree $0$, equal to
\begin{equation}
\label{eq:6:ZS-double-dual}
\mathbb{Z}[S_i]^{\solid}=\iHom\bigl(\Cont(S_i,\mathbb{Z}),\mathbb{Z}\bigr)
=\Bigl(\Cont(S_i,\mathbb{Z})\Bigr)^{\vee},
\end{equation}
the $\mathbb{Z}$-valued measures on $S_i$ (\cref{rmk:5:measures}). Each such term
is its own double dual, so $\mathbb{Z}[\underline{X}]^{\mathbb{L}\solid}$ is
computed by the complex of measures, which one rewrites as
\begin{equation}
\label{eq:6:solid-as-double-dual}
\mathbb{Z}[\underline{X}]^{\mathbb{L}\solid}\ \xrightarrow{\ \sim\ }\
R\iHom\Bigl(R\iHom\bigl(\mathbb{Z}[\underline{X}],\mathbb{Z}\bigr),\ \mathbb{Z}\Bigr).
\end{equation}
The inner $R\iHom(\mathbb{Z}[\underline{X}],\mathbb{Z})$ is the singular cochain
complex $C^{\bullet}_{\mathrm{sing}}(X,\mathbb{Z})$ by
\cref{thm:4:singular-recovery}; since $X$ is finite CW, each
$H^{i}_{\mathrm{sing}}(X,\mathbb{Z})$ is finitely generated, so dualizing back
turns cohomology into homology and yields
$C^{\mathrm{sing}}_{\bullet}(X,\mathbb{Z})$.
\end{proof}

\subsection{The structure of \texorpdfstring{$\Solid$}{Solid}}

We now describe the abelian category $\Solid$ more precisely. Since
$\prod_{\mathbb{N}}\mathbb{Z}$ is a compact projective generator, there are the
usual finiteness notions.

\begin{definition}[Finitely generated and finitely presented solid groups]
\label{def:6:fg-fp}
An object $M\in\Solid$ is \emph{finitely generated} if it is a quotient of a
single $\prod_{\mathbb{N}}\mathbb{Z}$, and \emph{finitely presented} if it is the
cokernel of a map
\begin{equation}
\label{eq:6:fp}
\prod_{\mathbb{N}}\mathbb{Z}\longrightarrow\prod_{\mathbb{N}}\mathbb{Z}.
\end{equation}
\end{definition}

Solid $\mathbb{Z}$-modules turn out to behave like modules over a coherent ring.

\begin{theorem}[Finitely presented solid groups form an abelian category]
\label{thm:6:fp-abelian}
The finitely presented objects of $\Solid$ form an abelian category, stable under
kernels, cokernels and extensions, and
\begin{equation}
\label{eq:6:ind-fp}
\Solid\ \simeq\ \mathrm{Ind}\bigl(\Solid^{\mathrm{fp}}\bigr).
\end{equation}
The only nontrivial closure property is stability under kernels; cokernels and
extensions then follow formally.
\end{theorem}

\begin{remark}[Resolutions of length one]
\label{rmk:6:length-one}
This is slightly stronger than what was announced in \cref{lec:1}: Clausen stated
that finitely presented objects admit resolutions of length two, but in fact
length one suffices. Every finitely presented $M\in\Solid$ has a resolution
\begin{equation}
\label{eq:6:fp-resolution}
0\longrightarrow\prod_{\mathbb{N}}\mathbb{Z}\longrightarrow\prod_{\mathbb{N}}\mathbb{Z}
\longrightarrow M\longrightarrow0,
\end{equation}
with the first map \emph{injective}; equivalently, finitely presented objects are
exactly the cokernels of injective maps of products of copies of $\mathbb{Z}$.
\end{remark}

The engine of both statements is the following.

\begin{lemma}[Key Lemma]
\label{lem:6:key}
Every finitely generated submodule of $\prod_{\mathbb{N}}\mathbb{Z}$ is isomorphic
to a product of copies of $\mathbb{Z}$ (an infinite product generically, possibly
zero or finite).
\end{lemma}

Granting \cref{lem:6:key}, \cref{thm:6:fp-abelian} follows: to see that the
finitely presented objects are stable under kernels, given a map with target $M$
one takes the kernel, which is a finitely generated submodule of a product, hence
by the Lemma itself a product $\prod_{\mathbb{N}}\mathbb{Z}$; the resolution
\eqref{eq:6:fp-resolution} results, and the rest is formal. The proof of
\cref{lem:6:key} uses countability and is genuinely a \emph{light} phenomenon
(the analogous statement fails for all condensed abelian groups). Its
combinatorial core is a classical fact of abelian group theory.

\begin{lemma}[Freeness fact]
\label{lem:6:free-fact}
Let $N$ be a countable abelian group that embeds into a product of copies of
$\mathbb{Z}$. Then $N$ is free.
\end{lemma}

\begin{remark}
\label{rmk:6:baer-specker}
Scholze remarked he was not sure whether this is stated in the literature, and
noted that the countability hypothesis is essential: an uncountable product
$\prod_{I}\mathbb{Z}$ is not free as an abstract abelian group.
\emph{Editorially:} this is the $\aleph_{1}$-freeness of the Baer--Specker group
$\prod_{\mathbb{N}}\mathbb{Z}$ --- every countable subgroup is free --- a classical
result of Specker; we record it only as a pointer.
\end{remark}

\begin{proof}[Proof of \cref{lem:6:key}]
Let $M\subseteq\prod_{\mathbb{N}}\mathbb{Z}$ be finitely generated, so that there
is a surjection $\prod_{\mathbb{N}}\mathbb{Z}\twoheadrightarrow M$ and an
injection $M\hookrightarrow\prod_{\mathbb{N}}\mathbb{Z}$; the composite
$g\colon\prod_{\mathbb{N}}\mathbb{Z}\to\prod_{\mathbb{N}}\mathbb{Z}$ is dual to a
map of direct sums
\begin{equation}
\label{eq:6:h-map}
h\colon\bigoplus_{\mathbb{N}}\mathbb{Z}\longrightarrow\bigoplus_{\mathbb{N}}\mathbb{Z},
\end{equation}
and the task is to show that $\operatorname{im}(g)$ is a product of copies of
$\mathbb{Z}$.

First reduce to $h$ injective. By \cref{lem:6:free-fact} the image
$\operatorname{im}(h)$, being countable and embedded in a direct sum (hence in a
product) of copies of $\mathbb{Z}$, is free; therefore
$\ker(h)\hookrightarrow\bigoplus_{\mathbb{N}}\mathbb{Z}$ splits off as a summand,
and after discarding it we may assume $h$ is injective, sitting in a short exact
sequence
\begin{equation}
\label{eq:6:h-ses}
0\longrightarrow\bigoplus_{\mathbb{N}}\mathbb{Z}\xrightarrow{\ h\ }
\bigoplus_{\mathbb{N}}\mathbb{Z}\longrightarrow Q\longrightarrow0.
\end{equation}
Next simplify $Q$. Let $\overline{Q}$ be the image of $Q$ under the natural map
into $\prod_{\varphi\in\Homop(Q,\mathbb{Z})}\mathbb{Z}$ (evaluate against all maps
$Q\to\mathbb{Z}$), so that there is a short exact sequence
\begin{equation}
\label{eq:6:Qbar}
0\longrightarrow Q'\longrightarrow Q\longrightarrow\overline{Q}\longrightarrow0,
\qquad \Homop(Q',\mathbb{Z})=0 .
\end{equation}
By \cref{lem:6:free-fact} $\overline{Q}$ is free, hence projective, so this
sequence splits and one may split $\overline{Q}$ off from
$\bigoplus_{\mathbb{N}}\mathbb{Z}$ as well; without loss of generality
\begin{equation}
\label{eq:6:Q-no-maps}
\Homop(Q,\mathbb{Z})=0 .
\end{equation}
(Concretely, $\overline{Q}$ is the cokernel of $Q$ in the category of countably
generated free abelian groups.)

Now dualize \eqref{eq:6:h-ses}. It gives an exact sequence
\begin{equation}
\label{eq:6:dualize}
0\to\Homop(Q,\mathbb{Z})\to\prod_{\mathbb{N}}\mathbb{Z}\xrightarrow{\ g=h^{\vee}\ }
\prod_{\mathbb{N}}\mathbb{Z}\to\Ext^{1}(Q,\mathbb{Z})\to0 .
\end{equation}
The point is that $\Homop(Q,\mathbb{Z})$ vanishes as a \emph{condensed} abelian
group, not merely on points: its $S$-valued points are
\begin{equation}
\label{eq:6:S-points}
\Homop\bigl(Q\otimes\mathbb{Z}[S],\mathbb{Z}\bigr)
=\Homop\bigl(Q,\Cont(S,\mathbb{Z})\bigr)=0,
\end{equation}
using tensor--Hom adjunction and the fact that $\Cont(S,\mathbb{Z})$ is a discrete
free abelian group, into which $Q$ (with no maps to $\mathbb{Z}$, hence to any
free abelian group) has no maps. Therefore $g$ is injective, and
$\operatorname{im}(g)\cong\prod_{\mathbb{N}}\mathbb{Z}$. The reductions made at the
start amounted precisely to arranging that $g$ became injective.
\end{proof}

\begin{corollary}[Classification of finitely presented solid groups]
\label{cor:6:fp-classification}
Every finitely presented $M\in\Solid$ is a direct sum of a product of copies of
$\mathbb{Z}$ and a group of the form
\begin{equation}
\label{eq:6:fp-form}
\Ext^{1}(Q,\mathbb{Z}),
\end{equation}
for some countable abelian group $Q$ with $\Homop(Q,\mathbb{Z})=0$. If $Q$ is
torsion then $\Ext^{1}(Q,\mathbb{Z})$ is profinite; in general one gets certain
(non-quasiseparated) solid groups.
\end{corollary}

\begin{corollary}[Flatness of $\prod_{\mathbb{N}}\mathbb{Z}$]
\label{cor:6:flat}
$\prod_{\mathbb{N}}\mathbb{Z}$ is flat for $\otimes^{\solid}$; and for solid $M$,
\begin{equation}
\label{eq:6:flat-formula}
M\otimes^{\solid}\prod_{\mathbb{N}}\mathbb{Z}\ \cong\ \prod_{\mathbb{N}}M .
\end{equation}
\end{corollary}

\begin{proof}
One must show $M\otimes^{\mathbb{L}\solid}\prod_{\mathbb{N}}\mathbb{Z}$ sits in
degree $0$ for every solid $M$; reducing to $M$ finitely presented, tensor the
resolution \eqref{eq:6:fp-resolution} with $\prod_{\mathbb{N}}\mathbb{Z}$. The
injective map stays injective, so the derived tensor product is concentrated in
degree $0$ and equals $\prod_{\mathbb{N}}M$.
\end{proof}

\begin{remark}[Flatness is special to the light setting]
\label{rmk:6:flat-light}
Abstractly a compact projective object need not be flat.
\Cref{cor:6:flat} in fact \emph{fails} for all (non-light) solid abelian groups,
but becomes true again for light condensed abelian groups --- which is why no
obstruction to it was ever seen in mathematical practice. Scholze recalled that
after passing to $p$-adic coefficients something like it was known, and that he
and Clausen had been puzzled about the solid case before observing this
light/heavy discrepancy.
\end{remark}

\subsection{Some \texorpdfstring{$\otimes^{\solid}$-}{solid tensor product }computations}

A striking feature of the solid tensor product is that many computations ``come
out right'', often in non-obvious ways, just as the earlier $\Ext$-computations
did. The relevant setting is $p$-adic formal geometry, where one takes derived
$p$-completions.

\begin{notation}[Derived $p$-completion]
\label{not:6:derived-p-completion}
For an abelian group (or complex) $M$, write
\begin{equation}
\label{eq:6:derived-p-completion}
M^{\wedge}_{p}:=R\varprojlim_{n}\ M/^{\mathbb{L}}p^{n},
\qquad
M/^{\mathbb{L}}p^{n}=\bigl[\,M\xrightarrow{p^{n}}M\,\bigr]\ (\text{in degrees }1,0),
\end{equation}
the \emph{derived} $p$-completion. Being derived $p$-complete can be checked on
cohomology groups. In the classical situation one often uses the completed tensor
product $\widehat{\otimes}$, the $p$-adic completion of the usual tensor product.
Almost always $M^{\wedge}_{p}$ agrees with the naive $p$-adic completion (it is in
degree $0$), but in full generality the derived version is the better operation.
\end{notation}

\begin{proposition}[$\otimes^{\mathbb{L}\solid}$ preserves derived $p$-completeness]
\label{prop:6:tensor-p-complete}
If $N,M\in D_{\geq0}(\Solid)$ are derived $p$-complete, then
$N\otimes^{\mathbb{L}\solid}M$ is derived $p$-complete.
\end{proposition}

\begin{corollary}
\label{cor:6:completed-sum-tensor}
For the completed free modules,
\begin{equation}
\label{eq:6:completed-sum-tensor}
\Bigl(\bigoplus_{\mathbb{N}}\mathbb{Z}\Bigr)^{\wedge}_{p}
\ \otimes^{\mathbb{L}\solid}\
\Bigl(\bigoplus_{\mathbb{N}}\mathbb{Z}\Bigr)^{\wedge}_{p}
\ \cong\
\Bigl(\bigoplus_{\mathbb{N}\times\mathbb{N}}\mathbb{Z}\Bigr)^{\wedge}_{p} .
\end{equation}
So in situations of formal geometry, where one completes, the tensor product does
what one expects.
\end{corollary}

\begin{remark}[Nothing special about $p$]
\label{rmk:6:general-ring}
There is nothing special about the number $p$ or about $\mathbb{Z}$. For any ring
$A$ and any element $x\in A$, let $\Solid_{A}$ denote the $A$-modules inside solid
abelian groups. Then for $M,N\in D_{\geq0}(\Solid_{A})$ derived $x$-complete,
$M\otimes^{\mathbb{L}\solid}_{A}N$ is again derived $x$-complete. (Here
$\Solid_{A}$ means, for now, simply $A$-modules in solid abelian groups.)
\end{remark}

The proof of \cref{prop:6:tensor-p-complete}, in the model case, exhibits the
combinatorial mechanism behind all such statements.

\begin{proof}[Proof sketch of \cref{cor:6:completed-sum-tensor}]
First, one may work with $\mathbb{Z}_p$ in place of $\mathbb{Z}$, because
$\mathbb{Z}_p\otimes^{\mathbb{L}\solid}\mathbb{Z}_p=\mathbb{Z}_p$: indeed
\begin{equation}
\label{eq:6:Zp-tensor}
\mathbb{Z}[\![T]\!]\otimes^{\mathbb{L}\solid}\mathbb{Z}[\![U]\!]
=\mathbb{Z}[\![U,T]\!]
\end{equation}
by \cref{cor:5:tensor-square}, and quotienting by $(T-p,\,U-p)$ gives
$\mathbb{Z}_p$ on each factor while quotienting the two-variable power series ring
by $T=p$ yields $\mathbb{Z}_p[\![U]\!]$, then further by $U=p$ gives $\mathbb{Z}_p$.
Consequently
\begin{equation}
\label{eq:6:D-solid-Zp}
D(\Solid_{\mathbb{Z}_p})\ \subset\ D(\Solid)
\end{equation}
is a full subcategory: being a $\mathbb{Z}_p$-module is a \emph{condition}, not a
structure, on a solid abelian group, and anything derived $p$-complete acquires
it. Take $M,N$ in this subcategory.

Consider the case $N=M=\bigl(\bigoplus_{\mathbb{N}}\mathbb{Z}_p\bigr)^{\wedge}_{p}$.
Writing it in terms of the compact projective generators is a small exercise: one
has
\begin{equation}
\label{eq:6:completed-sum-colim}
\Bigl(\bigoplus_{\mathbb{N}}\mathbb{Z}_p\Bigr)^{\wedge}_{p}\ \cong\
\operatorname*{colim}_{\substack{f\colon\mathbb{N}\to\mathbb{N}\\ f\to\infty}}\
\prod_{n\in\mathbb{N}}p^{f(n)}\mathbb{Z}_p ,
\end{equation}
the colimit over functions $f$ tending to $\infty$ (each fixed at value $0$ or
$1$ or\ldots only finitely often). There is an obvious map from each
$\prod_{n}p^{f(n)}\mathbb{Z}_p$ into the completed direct sum
$\bigl(\bigoplus_{\mathbb{N}}\mathbb{Z}_p\bigr)^{\wedge}_{p}
=\varprojlim_{n}\bigl(\bigoplus_{\mathbb{N}}\mathbb{Z}/p^{n}\bigr)$, since
$f\to\infty$ forces almost all coordinates to vanish modulo any fixed power of
$p$. It is injective; surjectivity is the point. Given
$S\in\ProFin_{\mathbb{N}}(\mathrm{Fin})$ and a map
$g\colon S\to\bigl(\bigoplus_{\mathbb{N}}\mathbb{Z}_p\bigr)^{\wedge}_{p}
=\varprojlim_{n}\bigl(\bigoplus_{\mathbb{N}}\mathbb{Z}/p^{n}\bigr)$ with components
$g_n$, each $g_n$ factors --- $S$ being compact and the target a direct sum ---
through $\bigoplus_{n\leq a(n)}\mathbb{Z}/p^{n}$ for some $a(n)$; choosing a
sufficiently slowly increasing $f$ makes $g$ factor through
$\prod_{n}p^{f(n)}\mathbb{Z}_p$.

Now compute the tensor product, pulling the colimits out (which
$\otimes^{\solid}$ commutes with):
\begin{align}
\label{eq:6:tensor-colim}
M\otimes^{\mathbb{L}\solid}N
&=\Bigl(\operatorname*{colim}_{f\to\infty}\prod_{n}p^{f(n)}\mathbb{Z}_p\Bigr)
\otimes^{\mathbb{L}\solid}
\Bigl(\operatorname*{colim}_{g\to\infty}\prod_{m}p^{g(m)}\mathbb{Z}_p\Bigr) \notag\\
&=\operatorname*{colim}_{f,g\to\infty}\ \prod_{n,m\in\mathbb{N}}p^{f(n)+g(m)}\mathbb{Z}_p
\ \xrightarrow{\ \sim\ }\
\operatorname*{colim}_{\substack{h\colon\mathbb{N}\times\mathbb{N}\to\mathbb{N}\\ h\to\infty}}
\prod_{n,m}p^{h(n,m)}\mathbb{Z}_p
=\Bigl(\bigoplus_{\mathbb{N}\times\mathbb{N}}\mathbb{Z}_p\Bigr)^{\wedge}_{p},
\end{align}
using that products of $\mathbb{Z}_p$'s tensor to products (the two-variable form
of \eqref{eq:6:generator-tensor-recall}). The final isomorphism is a
combinatorial lemma.
\end{proof}

\begin{lemma}[Combinatorial domination]
\label{lem:6:combinatorial}
Any function $h\colon\mathbb{N}\times\mathbb{N}\to\mathbb{N}$ tending to $\infty$
dominates an (even more slowly increasing) function, still tending to $\infty$,
that is a sum $f(n)+g(m)$ of functions of the individual coordinates. (In the
dual situation of \cref{prop:6:frechet} one uses the same fact with the
inequality reversed.)
\end{lemma}

At first sight the colimit over sums $f(n)+g(m)$ looks far smaller than the
colimit over all two-variable functions $h$; the lemma says they are cofinal.
Scholze left the (not hard) proof as an exercise. The general case of
\cref{prop:6:tensor-p-complete} reduces $M,N$ to completed direct sums of
generators and runs the same argument, with an extra step reducing an uncountable
index set to the countable case (derived completions commute with the relevant
colimits).

\subsection{Solid functional analysis}

Further computations come out right when one does $p$-adic functional analysis
inside the solid framework. Work over $\mathbb{Q}_p$; most of what follows works
over any complete nonarchimedean field. Inverting $p$ and passing to
$\mathbb{Z}_p$-modules gives a chain of full subcategories
\begin{equation}
\label{eq:6:Qp-subcats}
D(\Solid_{\mathbb{Q}_p})\ \subset\ D(\Solid_{\mathbb{Z}_p})\ \subset\ D(\Solid).
\end{equation}
Tensor products stay within each subcategory. The category $\Solid_{\mathbb{Q}_p}$
of solid $\mathbb{Q}_p$-vector spaces has a compact projective generator
\begin{equation}
\label{eq:6:smith-generator}
\Bigl(\prod_{\mathbb{N}}\mathbb{Z}_p\Bigr)\bigl[\tfrac1p\bigr],
\end{equation}
a $p$-adic \emph{Smith space}, while the more usual objects are the $p$-adic
Banach spaces
\begin{equation}
\label{eq:6:banach}
\Bigl(\bigoplus_{\mathbb{N}}\mathbb{Z}_p\Bigr)^{\wedge}_{p}\bigl[\tfrac1p\bigr].
\end{equation}

\begin{remark}[Smith spaces are not normed]
\label{rmk:6:smith-not-normed}
The generator \eqref{eq:6:smith-generator} is a curious condensed vector space
that does come from a topological vector space, but is not a normed space. It
looks like one with a distinguished ``unit ball'' $\prod_{\mathbb{N}}\mathbb{Z}_p$
that is \emph{compact}; but no infinite-dimensional Banach (or normed) space has
compact unit ball, so if one tried to endow it with a norm making that set the
unit ball, the norm would fail to be continuous. Such objects have appeared in
classical functional analysis under the name \emph{Smith spaces}: one considers
topological vector spaces that are increasing unions obtained by scaling out a
fixed compact convex set, dually to Banach spaces. Scholze attributed the notion
to Marianne F.\ Smith, who studied the analogue over $\mathbb{R}$.
\end{remark}

The Smith spaces and (separable) Banach spaces are exchanged by duality.

\begin{proposition}[Smith--Banach duality]
\label{prop:6:smith-banach}
There is an anti-equivalence
\begin{equation}
\label{eq:6:smith-banach}
\bigl(\{\text{light Smith spaces}\}\bigr)^{\mathrm{op}}\ \simeq\
\{\text{separable Banach spaces}\},
\end{equation}
given in both directions by $\mathbb{Q}_p$-linear duality
$V\mapsto\iHom(V,\mathbb{Q}_p)$, $W\mapsto\iHom(W,\mathbb{Q}_p)$.
\end{proposition}

Concretely, the dual of the Smith space $(\prod_{\mathbb{N}}\mathbb{Z}_p)[\tfrac1p]$
is the Banach space $(\bigoplus_{\mathbb{N}}\mathbb{Z}_p)^{\wedge}_p[\tfrac1p]$
and vice versa; ``separable/light'' on one side matches a countable direct sum,
on the other a countable product.

\begin{remark}[Derived duality and set theory]
\label{rmk:6:derived-duality}
One may ask what happens to duality \emph{derived}. Going from a Smith space $V$,
\begin{equation}
\label{eq:6:derived-dual-smith}
R\iHom(V,\mathbb{Q}_p)=\iHom(V,\mathbb{Q}_p),
\end{equation}
with no higher terms --- true and essentially obvious, since $V$ is internally
projective. Going the other way, from a Banach space $W$, the question whether
\begin{equation}
\label{eq:6:derived-dual-banach}
\iHom(W,\mathbb{Q}_p)\ \overset{?}{=}\ R\iHom(W,\mathbb{Q}_p)
\end{equation}
\emph{depends on the model of set theory}, exactly as with the phenomenon of
\cref{lec:4}: it fails under the continuum hypothesis, but holds under the
principle $(\ast)$ of \cref{def:4:star}, which is consistent
(\cref{thm:4:star-set-theory}). The reason is that a Banach space is a special
pro-$p$-type object, a limit of direct sums
$\varprojlim_n\bigl(\bigoplus_{\mathbb{N}}\mathbb{Q}_p/p^{n}\mathbb{Z}_p\bigr)$, so
the higher $\Ext$ against $\mathbb{Q}_p$ is precisely controlled by $(\ast)$.
Scholze said he was increasingly tempted to just work in a model where $(\ast)$
holds, so that these higher $\Ext$'s are $0$ rather than meaningless junk; in the
computations of solid functional analysis (many duals and tensor products) one
does at some point take an $R\iHom$ of a Banach space against $\mathbb{Q}_p$, and
then one may choose the well-behaved model. He and Clausen never found the
question genuinely relevant, but it does come up.
\end{remark}

Finally, the solid tensor product recovers the classical completed tensor
products of functional analysis --- even though it was defined to be compatible
with \emph{colimits}, whereas those tensor products are usually designed to be
compatible with \emph{limits}.

\begin{proposition}[Fréchet spaces]
\label{prop:6:frechet}
A \emph{Fréchet space} is a countable limit of Banach spaces along dense
transition maps; these carry a standard completed tensor product
$\widehat{\otimes}$, compatible with such limits. For Fréchet
$\mathbb{Q}_p$-vector spaces $V,W$,
\begin{equation}
\label{eq:6:frechet-tensor}
V\otimes^{\mathbb{L}\solid}W\ \cong\ V\,\widehat{\otimes}\,W,
\end{equation}
concentrated in degree $0$. For example
\begin{equation}
\label{eq:6:prod-Qp-tensor}
\prod_{\mathbb{N}}\mathbb{Q}_p\ \otimes^{\mathbb{L}\solid}\ \prod_{\mathbb{N}}\mathbb{Q}_p
\ \cong\ \prod_{\mathbb{N}\times\mathbb{N}}\mathbb{Q}_p .
\end{equation}
\end{proposition}

\begin{proof}[Proof in the example]
Write the Fréchet space $\prod_{\mathbb{N}}\mathbb{Q}_p$ as a colimit of Smith
spaces with ever larger denominators,
\begin{equation}
\label{eq:6:prod-Qp-colim}
\prod_{\mathbb{N}}\mathbb{Q}_p\ =\
\operatorname*{colim}_{\substack{f\colon\mathbb{N}\to\mathbb{N}\\ f\to\infty\ \text{fast}}}\
\Bigl(\prod_{n}p^{-f(n)}\mathbb{Z}_p\Bigr)\bigl[\tfrac1p\bigr] ,
\end{equation}
the colimit now over \emph{very fast} increasing $f$ (any profinite set mapping
in factors through some $\prod_n p^{-f(n)}\mathbb{Z}_p$). Then
\begin{equation}
\label{eq:6:prod-Qp-tensor-colim}
\prod_{\mathbb{N}}\mathbb{Q}_p\otimes^{\mathbb{L}\solid}\prod_{\mathbb{N}}\mathbb{Q}_p
=\operatorname*{colim}_{f,g\to\infty}\Bigl(\prod_{\mathbb{N}\times\mathbb{N}}
p^{-f(n)-g(m)}\mathbb{Z}_p\Bigr)\bigl[\tfrac1p\bigr]
=\prod_{\mathbb{N}\times\mathbb{N}}\mathbb{Q}_p,
\end{equation}
again by the combinatorial domination \cref{lem:6:combinatorial}.
\end{proof}

\begin{remark}[A huge space, not a compact projective one]
\label{rmk:6:huge-not-compact}
The answer $\prod_{\mathbb{N}\times\mathbb{N}}\mathbb{Q}_p$ in
\eqref{eq:6:prod-Qp-tensor} might look like a formal statement about compact
projective generators, but it is not: the spaces $\prod_{\mathbb{N}}\mathbb{Q}_p$
are very far from compact projective (only the unit balls
$(\prod_{\mathbb{N}}\mathbb{Z}_p)[\tfrac1p]$ are). The general proof combines the
argument for the compact case with the present one; Scholze noted Efimov has a
slicker way to combine them, used in his categorical work.
\end{remark}

\subsection{Questions}

\begin{remark}[Light vs.\ all condensed sets for Banach spaces]
\label{rmk:6:light-banach}
For functional analysis in large Banach spaces it makes no difference whether one
works in light or all condensed sets: all Banach spaces --- and all Fréchet
spaces --- are light, separable or not, since even an uncountable-dimensional
Banach space is a naive filtered colimit of countable-dimensional (separable)
Banach subspaces. It is only the predual Smith spaces that see the distinction.
\end{remark}

\begin{remark}[Solidification of $K$-theory spectra]
\label{rmk:6:K-theory}
The relation between the singular complex and solidification
(\cref{prop:6:singular-homology}) is no accident: one can replace $\CondAbl$
everywhere by condensed spectra and solidify there, since $\Solid$ sits inside
all condensed spectra just as solid abelian groups sit inside condensed abelian
groups. A telling example is connective complex $K$-theory $ku$. The algebraic
$K$-theory spectrum $K(\mathbb{C})$ is naturally a \emph{condensed} spectrum:
$\mathbb{C}$ is a condensed ring, and one sets
$S\mapsto K\bigl(\Cont(S,\mathbb{C})\bigr)$. This is not yet the topological
object --- e.g.\ $\pi_1$ is $\mathbb{C}^{\times}$ as a condensed abelian group ---
but its solidification contracts the groups $\mathrm{GL}_n(\mathbb{C})$ down to
their homotopy types (by the mechanism of \cref{prop:6:singular-homology}), and
one recovers topological connective $ku$. More generally, for any
$\mathbb{C}$-algebra $A$ (with any topology, even discrete) one may form the
condensed spectrum $K(A)$ and solidify, obtaining the \emph{semi-topological}
$K$-theory of $A$; the analogous constructions over $\mathbb{R}$ give $KO$, etc.
The output is connective (not entirely obviously); recovering the periodic $KU$
--- inverting the Bott element --- is more subtle, and is one motivation for the
\emph{analytic} $K$-theory Clausen and Scholze developed, which allows negative
degrees. The underlying mechanism is that solidification factors through passage
to homotopy types on suitable objects: in the ambient category of condensed
anima, the ``homotopy type'' functor is a left adjoint defined on a class of
objects (e.g.\ on classifying spaces $B\mathrm{GL}_n(\mathbb{C})$), and
solidification applies it.
\end{remark}
